\documentclass[11pt,a4paper]{article}
\usepackage{stile}

\begin{document}
\documentotitolo{Manuale utente}{Alberi di regressione --- esecuzione dei JAR}

Il manuale descrive l'avvio dei JAR tramite script, soprattutto su Windows,
e l'uso del client per apprendere alberi, caricare archivi e ottenere predizioni.
È disponibile anche un'applicazione locale su file, senza MySQL né server.

\section*{Prerequisiti}

Estrarre tutto il progetto in una cartella scrivibile, per esempio Documenti.
Non avviare i file dallo ZIP e mantenere la struttura di \texttt{distribution}.

\subsection*{Java su Windows}

Installare \textbf{Java 8 o successivo}, se non è già disponibile.
Gli script eseguono i JAR pronti: \textbf{non compilano il progetto e non
richiedono Eclipse o Maven}. Cercano Java prima in \texttt{JAVA\_HOME} e,
se questa variabile non è impostata, nel \texttt{PATH}.

Per configurare Java, cercare nel menu Start \textbf{Modifica le variabili di
ambiente per l'account}. Creare la variabile utente \texttt{JAVA\_HOME} con
il percorso della cartella del JDK, \textbf{senza virgolette e senza}
\texttt{\textbackslash bin}. Aprire un nuovo \textbf{Prompt dei comandi} e verificare:
\begin{lstlisting}
"%JAVA_HOME%\bin\java.exe" -version
\end{lstlisting}

In alternativa, senza \texttt{JAVA\_HOME}, deve funzionare \texttt{java -version}.
Se \texttt{JAVA\_HOME} è errata, correggerla oppure rimuoverla per usare il
\texttt{PATH}. Gli script non usano la configurazione Java di Eclipse.

\subsection*{MySQL: solo per l'apprendimento client/server}

MySQL deve essere installato e in esecuzione sulla stessa macchina del server,
sulla porta \texttt{3306}. Prima del primo utilizzo, aprire in
\textbf{MySQL Workbench} una connessione locale con un account amministrativo;
con \textbf{File > Open SQL Script} aprire ed eseguire l'intero file:
\begin{lstlisting}
distribution/server/sql_file/map6_setup.sql
\end{lstlisting}

Il file crea \texttt{MapDB}, le tabelle \texttt{provaC} e \texttt{servo} e
l'utente \texttt{MapUser} con password \texttt{map}. Verificare che termini
senza errori. È per la \textbf{prima installazione}: non rieseguirlo su un
database già configurato senza controllarne il contenuto.
Gli script di avvio non installano né configurano MySQL.

\clearpage
\section*{Avvio su Windows}

Aprire la cartella \texttt{distribution} in Esplora file. I tre script sono:
\begin{itemize}
  \item \texttt{avvia-server.bat}: esegue \texttt{server/server.jar};
  \item \texttt{avvia-client.bat}: esegue \texttt{client/client\_base.jar};
  \item \texttt{avvia-locale.bat}: esegue \texttt{src/src.jar}, senza server.
\end{itemize}

\subsection*{Server}

Dopo aver preparato MySQL, fare doppio clic su \texttt{avvia-server.bat}.
Senza argomenti viene utilizzata la porta \textbf{8080}.
Lasciare aperta la finestra e avviare il client come descritto nella pagina
seguente. Non fare doppio clic direttamente sul JAR: l'applicazione richiede
una console e gli script impostano la cartella di lavoro corretta.

Lo script stampa un messaggio prima di eseguire Java; non è una verifica
che il servizio sia già pronto. Il server non stampa una conferma quando
inizia ad attendere connessioni: se il processo rimane attivo senza altri
messaggi, si può provare a collegare il client. Se invece compare un errore
o il messaggio \texttt{Esecuzione terminata}, il server non è in ascolto.

Gli alberi appresi vengono salvati in \texttt{distribution/server}, anche
se lo script è stato richiamato da un'altra cartella. Il driver
\texttt{mysql-connector-java-8.0.17.jar} deve rimanere accanto a
\texttt{server.jar}.

\subsection*{Porta e indirizzo personalizzati}

Il doppio clic usa \texttt{localhost:8080}, adatto a server e client sullo
stesso computer. Per cambiare porta, aprire due Prompt dei comandi nella
cartella principale del progetto ed eseguire, prima nel primo e poi nel secondo:
\begin{lstlisting}
distribution\avvia-server.bat 9090
\end{lstlisting}
\begin{lstlisting}
distribution\avvia-client.bat localhost 9090
\end{lstlisting}

Se il server è su un altro computer, sostituire \texttt{localhost} con il suo
indirizzo IP e consentire la porta TCP del server nel firewall, sulla rete
utilizzata. La porta deve coincidere nei due script; usare un numero tra
\texttt{1} e \texttt{65535}. Non confondere questa porta con la \texttt{3306}
di MySQL. In PowerShell anteporre \texttt{.\textbackslash} ai percorsi degli
script, per esempio \texttt{.\textbackslash distribution\textbackslash avvia-server.bat}.

\subsection*{Arresto}

Per arrestare il server premere \textbf{Ctrl+C} nella sua finestra e confermare
l'eventuale richiesta di terminare il processo batch. Gli script Windows
attendono la pressione di un tasto al termine dell'applicazione, anche in caso
di errore, per permettere di leggere i messaggi.

\clearpage
\section*{CLI}

La versione CLI utilizza il terminale per mostrare il menu, richiedere i dati e
visualizzare il risultato. Dopo ogni scelta digitata occorre premere Invio.
Negli esempi seguenti le righe in \textcolor{MAPinput}{blu} indicano l'input
dell'utente.

\subsection*{Client}

Con il server già avviato sulla porta 8080, fare doppio clic su
\texttt{distribution/avvia-client.bat}. Si apre una \textbf{seconda finestra},
collegata a \texttt{localhost:8080}. In alternativa, dal Prompt dei comandi
nella cartella principale del progetto:

\begin{lstlisting}
distribution\avvia-client.bat
\end{lstlisting}

Per un indirizzo o una porta diversi, passare gli argomenti allo script come
indicato nella sezione precedente. Il server va sempre avviato prima del client.

Una volta stabilita la connessione, il client presenta il menu:

\begin{lstlisting}
Learn Regression Tree from data [1]
Load Regression Tree from archive [2]
\end{lstlisting}
\didascalia{Menu principale del client da console.}

\begin{itemize}
  \item \textbf{1}: acquisisce gli esempi da una tabella del database e apprende
        un nuovo albero;
  \item \textbf{2}: carica dal disco del server un albero salvato in una
        precedente esecuzione.
\end{itemize}

\subsection*{Apprendimento da database}

Selezionare \textbf{1}. Alla richiesta \texttt{Table name:}, inserire il nome
della tabella da utilizzare, per esempio \texttt{provaC}.

\begin{lstlisting}
Learn Regression Tree from data [1]
Load Regression Tree from archive [2]
(*@\inpututente{1}@*)
Table name:
(*@\inpututente{provaC}@*)
Table found!
Starting data acquisition phase!
Starting learning phase!
\end{lstlisting}
\didascalia{Selezione della tabella e avvio dell'apprendimento.}

Dopo il caricamento dei dati, il server costruisce l'albero e lo salva nel file
\texttt{provaC.dmp}. Se l'operazione termina correttamente, il client passa
alla fase di predizione.

\clearpage
\subsection*{Predizione}

Durante questa fase il client mostra i rami disponibili nel nodo corrente.
Ogni riga contiene un indice e la condizione associata: per scegliere il ramo
occorre digitare il suo \textbf{indice numerico}, non il nome dell'attributo
o il valore da predire.

Nell'esempio con \texttt{provaC}, il primo nodo distingue i valori
\texttt{A} e \texttt{B} dell'attributo \texttt{X}:

\begin{lstlisting}
Starting prediction phase!
0:X=A
1:X=B
(*@\inpututente{0}@*)
\end{lstlisting}
\didascalia{Scelta del ramo corrispondente a X=A.}

Se il ramo selezionato conduce a un altro nodo interno, viene presentata una
nuova domanda. Il numero di scelte dipende dal percorso seguito nell'albero.
Proseguendo nell'esempio:

\begin{lstlisting}
0:Y<=2.0
1:Y>2.0
(*@\inpututente{0}@*)
\end{lstlisting}
\didascalia{Seconda scelta: il ramo in cui Y è minore o uguale a 2.0.}

Raggiunta una foglia, il client mostra il valore predetto, preceduto dal
messaggio \texttt{Predicted class:}. Per le scelte indicate il risultato è
\texttt{1.0}.

\begin{lstlisting}
Predicted class:1.0
Would you repeat ? (y/n)
\end{lstlisting}
\didascalia{Risultato della predizione e richiesta di ripetizione.}

Alla domanda finale:
\begin{itemize}
  \item digitare \textbf{y} per iniziare una nuova predizione sullo stesso
        albero, ripartendo dalla radice;
  \item digitare \textbf{n} per terminare il client.
\end{itemize}

La chiusura del client non arresta il server, che resta disponibile per altre
connessioni.

\clearpage
\subsection*{Caricamento da archivio}

L'opzione \textbf{2} del menu iniziale permette di riutilizzare un albero già
salvato. Prima di selezionarla deve quindi essere disponibile un archivio
prodotto dal server, per esempio \texttt{provaC.dmp} dopo l'apprendimento
descritto nelle pagine precedenti.

Il file deve trovarsi nella directory di lavoro del server. Se il servizio è
stato avviato con gli script di questo manuale, tale directory è
\texttt{distribution/server}. L'archivio non va cercato nella cartella del
client.

Avviare nuovamente il client, selezionare \textbf{2} e rispondere alla richiesta
\texttt{File name:} inserendo il nome \textbf{senza l'estensione}
\texttt{.dmp}:

\begin{lstlisting}
Learn Regression Tree from data [1]
Load Regression Tree from archive [2]
(*@\inpututente{2}@*)
File name:
(*@\inpututente{provaC}@*)
Table found!
Starting prediction phase!
\end{lstlisting}
\didascalia{Caricamento dell'albero conservato in provaC.dmp.}

Il messaggio \texttt{Table found!} viene utilizzato dal programma anche per
confermare il caricamento di un archivio. Non viene eseguito un nuovo
apprendimento: la predizione utilizza l'albero letto dal file.

Le scelte si effettuano con la stessa procedura già descritta. Nell'esempio,
selezionando il ramo \texttt{1}, corrispondente a \texttt{X=B}, si raggiunge
una foglia con valore \texttt{10.0}:

\begin{lstlisting}
0:X=A
1:X=B
(*@\inpututente{1}@*)
Predicted class:10.0
Would you repeat ? (y/n)
(*@\inpututente{n}@*)
\end{lstlisting}
\didascalia{Predizione sull'albero caricato e chiusura del client.}

Anche dopo il caricamento è possibile rispondere \textbf{y} per ripetere la
predizione oppure \textbf{n} per terminare.

\clearpage
\section*{Applicazione locale: prova senza MySQL}

Per usare i dataset su file, senza server né client, aprire la cartella
\texttt{distribution} e fare doppio clic su \texttt{avvia-locale.bat}.
Dal Prompt dei comandi nella cartella principale del progetto:
\begin{lstlisting}
distribution\avvia-locale.bat
\end{lstlisting}

Lo script esegue \texttt{src/src.jar} con directory di lavoro
\texttt{distribution/src}, dove sono già presenti \texttt{prova.dat},
\texttt{provaC.dat} e \texttt{servo.dat}.

\begin{enumerate}
  \item Scegliere \textbf{1} per apprendere un albero e, alla richiesta
        \texttt{File name:}, digitare per esempio \texttt{provaC},
        \textbf{senza} \texttt{.dat}.
  \item Il programma legge il dataset, salva \texttt{provaC.dmp} nella stessa
        cartella e stampa le regole dell'albero.
  \item Per la predizione scegliere gli indici numerici dei rami proposti.
        Il valore predetto viene stampato a video; rispondere \textbf{y} per
        ripetere oppure \textbf{n} per terminare.
  \item A un avvio successivo, scegliere \textbf{2} e digitare
        \texttt{provaC}, senza \texttt{.dmp}, per caricare l'albero salvato.
\end{enumerate}

Gli archivi del server e dell'applicazione locale \textbf{non sono
intercambiabili}: mantenerli nelle rispettive cartelle. Anche i messaggi della
console locale possono differire da quelli del client illustrati in precedenza.

\section*{Eclipse e JAR distribuiti}

Per provare i JAR è sufficiente usare gli script da Esplora file o dal terminale,
anche quando Eclipse è installato. Non occorre importare i progetti.

Solo per consultare o modificare i sorgenti, in Eclipse scegliere
\textbf{File > Import > General > Existing Projects into Workspace}, selezionare
\texttt{project} e abilitare \textbf{Search for nested projects}. Importare
\texttt{mapClient}, \texttt{mapServer} e \texttt{src} senza selezionare
\textbf{Copy projects into workspace}, così da mantenere il collegamento al
driver JDBC in \texttt{distribution/server}.

Per compilare i sorgenti serve un JDK compatibile con la versione dell'IDE.
In \textbf{Window > Preferences > Java > Installed JREs} aggiungere il JDK;
in \textbf{Java > Installed JREs > Execution Environments} associarlo a
\texttt{JavaSE-1.8}. Questa impostazione riguarda Eclipse, non
\texttt{JAVA\_HOME} o il \texttt{PATH} usati dagli script.

\textbf{Run As > Java Application} esegue le classi compilate da Eclipse,
non i JAR distribuiti. Dopo una modifica ai sorgenti, gli script continuano a
eseguire i vecchi JAR finché questi non vengono esportati nuovamente.

\clearpage
\section*{Linux e macOS}

Sono disponibili gli equivalenti \texttt{.sh}, con gli stessi valori predefiniti,
argomenti e directory di lavoro. Serve Java 8 o successivo disponibile tramite
\texttt{JAVA\_HOME} oppure \texttt{PATH}; per apprendere dal database occorre
preparare MySQL come descritto nei prerequisiti.

Dalla cartella principale del progetto, avviare il server:
\begin{lstlisting}
sh distribution/avvia-server.sh
\end{lstlisting}
Poi, in un secondo terminale, avviare il client:
\begin{lstlisting}
sh distribution/avvia-client.sh
\end{lstlisting}
Per usare una porta diversa o un altro indirizzo:
\begin{lstlisting}
sh distribution/avvia-server.sh 9090
sh distribution/avvia-client.sh localhost 9090
\end{lstlisting}
I due comandi vanno eseguiti in terminali separati. Per la sola applicazione locale:
\begin{lstlisting}
sh distribution/avvia-locale.sh
\end{lstlisting}
Gli script \texttt{.sh} non attendono un tasto al termine. Per interrompere
l'esecuzione utilizzare \textbf{Ctrl+C}.

\section*{Problemi comuni}

\begin{description}[style=nextline,leftmargin=0pt]
  \item[Java non trovato o JAVA\_HOME non valida.]
    Verificare i passaggi dei prerequisiti e riaprire il terminale dopo aver
    modificato le variabili di ambiente. Un errore \texttt{UnsupportedClassVersionError}
    indica che occorre una versione di Java più recente, almeno Java 8 per i JAR forniti.
  \item[JAR o driver mancante.]
    Estrarre di nuovo l'intera distribuzione mantenendo cartelle e nomi originali;
    non spostare gli script separatamente dai JAR.
  \item[Connection refused oppure errore di connessione del client.]
    Verificare che il server sia ancora attivo e che indirizzo e porta coincidano.
    Per collegamenti tra computer verificare anche rete e firewall.
  \item[Address already in use nel server.]
    La porta è occupata: chiudere l'altra istanza del server oppure scegliere
    un'altra porta, passandola anche allo script del client.
  \item[Errore di database o tabella non trovata.]
    Controllare che MySQL sia attivo su \texttt{localhost:3306}, che lo script SQL
    sia stato completato e che l'utente \texttt{MapUser} possa accedere a
    \texttt{MapDB}. Le tabelle di esempio sono \texttt{provaC} e \texttt{servo}.
  \item[Archivio o dataset non trovato; salvataggio non riuscito.]
    Inserire i nomi senza estensione. Gli archivi del server sono in
    \texttt{distribution/server}; dataset e archivi locali sono in
    \texttt{distribution/src}. Per caricare un archivio occorre averlo prima
    generato tramite apprendimento. Verificare che la cartella sia scrivibile.
\end{description}

\end{document}
